\documentclass[11pt]{article}
\usepackage[margin=1in]{geometry}
\usepackage{amsmath,amssymb}
\title{Disproof of Conjecture 00000007623}
\author{earthking11}
\date{6 October 2026}
\begin{document}
\maketitle

The conjecture says that the nilpotency index is at most two, with a
``degenerate'' exception of index three.  Consider the linear operator on
$\mathbb Q^4$ whose matrix in the standard basis is
\[
 J=\begin{pmatrix}
 0&1&0&0\\
 0&0&1&0\\
 0&0&0&1\\
 0&0&0&0
 \end{pmatrix}.
\]
For a coordinate vector, $J(a,b,c,d)=(b,c,d,0)$.  Consequently
\[
 J^4=0,
 \qquad
 J^3(0,0,0,1)=(1,0,0,0)\ne0.
\]
Thus $J$ is nilpotent of index exactly four.  This is neither an index at
most two nor the stated index-three exception, so the conjecture is false.

The Lean project encodes the same four-coordinate operator over the integers.
The fourth-iterate identity is definitional, while the displayed nonzero
third iterate is kernel-checked.  No search or unformalized computation is
needed.
\end{document}
